\documentclass[a4paper,11pt,center]{article}
\usepackage[utf8]{inputenc}
\usepackage{graphicx}

\begin{document}


\begin{center}
	\textbf{MOUNTAINS OF THE MOON UNIVERSITY\newline\newline\newline
		FACULTY OF SCIENCE TECHNOLOGY\newline\newline\newline
	    AND INNOVATION\newline\newline\newline
		DEPARTMENT OF COMPUTER SCIENCE\newline\newline\newline
		Coursework}\newline\newline\newline
\end{center}
		

	
\textbf{NAME: MUHINDO JOSEPH }\newline\newline
	\textbf{REGISTRATION NUMBER: 230270/U/MMU/BCS/01139}\newline\newline
	\textbf{COURSE CODE: BCS 1201}\newline\newline
	\textbf{LECTURER'S NAME: Mr. Samuel Ocen}

\newpage


\section{Resource Allocation in Cloud 
	Computing}
\subsection{\textbf{Number One: Basic Resource Allocation}}


\begin{verbatim}
	#importing necessary libraries
	import numpy as np
	import matplotlib.pyplot as plt
	from scipy.optimize import linprog
	
	# Define the objective function coefficients
	c = [4, 5]
	
	# Coefficients of the inequality constraints
	A = np.array([
	[-2, -3],  # CPU constraint
	[-1, -2],  # Memory constraint
	[-3, -1],  # Storage constraint
	])
	
	# Right-hand side values of the inequality constraints
	b = [-10, -5, -8]
	
	# Solve linear programming problem
	result = linprog(c, A_ub=A, b_ub=b)
	
	# Display the results
	print("Optimal values:")
	print("x =", result.x[0])
	print("y =", result.x[1])
	print("Optimal objective function value (z) =", result.fun)
	
	# Plotting the feasible region
	x_values = np.linspace(0, 5, 100)
	y1_values = (10 - 2*x_values) / 3
	y2_values = (5 - x_values) / 2
	y3_values = 8 - 3*x_values
	
	plt.figure(figsize=(8, 6))
	plt.plot(x_values, y1_values, label='2x + 3y >= 10')
	plt.plot(x_values, y2_values, label='x + 2y >= 5')
	plt.plot(x_values, y3_values, label='3x + y >= 8')
	
	plt.fill_between(x_values, np.maximum(np.maximum(y1_values, y2_values), y3_values), color='green', alpha=0.5, label='Unwanted Region')
	plt.scatter(result.x[0], result.x[1], color='red', marker='*', label='Optimal Solution')
	
	plt.xlabel('x')
	plt.ylabel('y')
	plt.ylim(0,10)
	plt.xlim(0,6)
	plt.title('Basic Resource Allocation')
	plt.legend()
	plt.grid(True)
	plt.show()
	
\end{verbatim}

\textbf{Solution And Graph}\newline\newline
\includegraphics{solutions/num1}\newline\newline
\includegraphics{C:/Users/skurge/OneDrive/Desktop/graph1}
\newpage




\subsection{\textbf{Number Two: Load Balancing}}

\begin{verbatim}
	import numpy as np
	import matplotlib.pyplot as plt
	from scipy.optimize import linprog
	
	# Define the objective function coefficients
	c = [5, 4]
	
	# Coefficients of the inequality constraints
	A = np.array([
	[-2, -3],  # Server 1 constraint
	[-4, -2],  # Server 2 constraint
	])
	
	# Right-hand side values of the inequality constraints
	b = [-20, -15]
	
	# Solve linear programming problem
	result = linprog(c, A_ub=A, b_ub=b)
	
	# Display the results
	print("Optimal values:")
	print("x =", result.x[0])
	print("y =", result.x[1])
	print("Optimal objective function value (z) =", result.fun)
	
	# Plotting the feasible region
	x_values = np.linspace(0, 10, 100)
	y1_values = (20 - 2*x_values) / 3  # Server 1 constraint
	y2_values = (15 - 4*x_values) / 2  # Server 2 constraint
	
	plt.figure(figsize=(8, 6))
	plt.plot(x_values, y1_values, label='2x + 3y >= 20')
	plt.plot(x_values, y2_values, label='4x + 2y >= 15')
	
	
	x = [0.625, 0, 0, 10, 10]
	y = [6.25, 7.5, 10, 10, 0]
	plt.fill(x,y, color="green", label="Feasible Region")
	
	# plt.fill_between(x_values, np.maximum(y1_values, y2_values), color='gray', alpha=0.5, label='Feasible Region')
	plt.scatter(result.x[0], result.x[1], color='yellow', marker='*', label='Optimal Solution')
	
	plt.xlabel('x')
	plt.ylabel('y')
	plt.xlim(0, 10)
	plt.ylim(0, 9.0)
	plt.title('Load Balancing Problem')
	plt.legend()
	plt.grid(True)
	plt.show()
	
\end{verbatim}

\textbf{Solution And Graph}\newline\newline
\includegraphics{solutions/num2}\newline\newline
\includegraphics{C:/Users/skurge/OneDrive/Desktop/graph2}
\newpage




\subsection{\textbf{Number Three: Energy Efficient Resource Allocation}}


\begin{verbatim}
	from pulp import *
	import matplotlib.pyplot as plt
	import numpy as np
	
	model = LpProblem(name = "Energy_consumption", sense=LpMinimize)
	
	x = LpVariable("x", 0)
	y = LpVariable("y", 0)
	
	model += 3*x + 2*y
	
	model += 2*x + 3*y >= 15
	model += 4*x + 2*y >= 10
	
	model.solve()
	
	optimal_x = x.varValue
	optimal_y = y.varValue
	optimal_value = model.objective.value()
	
	print("Optimal Solution: ")
	print("x: ", optimal_x)
	print("y: ", optimal_y)
	print("Objective value: ", optimal_value)
	
	
	x_values = np.linspace(0, 10, 100)
	y1_values = (15-2*x_values)/3
	y2_values = (10-4*x_values)/2
	
	plt.plot(x_values, y1_values, label=2*x + 3*y >= 15)
	plt.plot(x_values, y2_values, label=4*x + 2*y >= 10)
	
	# plt.fill_between(x_values, np.maximum(y1_values,y2_values), 0, where=(y1_values<=y2_values), color="orange")
	plt.fill_between(x_values, np.maximum(y1_values,y2_values), 0, where=(y1_values>=y2_values), color="yellow", alpha=0.5, label='Unwanted Region')
	
	
	plt.title("Graph showing the feasible region")
	plt.xlim(0,3)
	plt.ylim(0,6)
	plt.grid(True)
	plt.scatter(optimal_x, optimal_y, color="red",marker='*', label="Optimal Solution")
	plt.legend()
	plt.show()
\end{verbatim}

\textbf{Solution And Graph}\newline\newline
\includegraphics{solutions/num3}\newline\newline
\includegraphics{C:/Users/skurge/OneDrive/Desktop/graph3}
\newpage




\subsection{\textbf{Number Four: Multi-Tenant Resource Sharing}}


\begin{verbatim}
	from pulp import *
	import matplotlib.pyplot as plt
	import numpy as np
	
	model = LpProblem(name = "Multi-Tenant-Resource-Sharing", sense=LpMinimize)
	
	x = LpVariable("x", 0)
	y = LpVariable("y", 0)
	
	model += 5*x + 4*y
	
	model += 2*x + 3*y >= 12
	model += 4*x + 2*y >= 18
	
	model.solve()
	
	optimal_x = x.varValue
	optimal_y = y.varValue
	optimal_value = model.objective.value()
	
	print("Optimal Solution: ")
	print("x: ", optimal_x)
	print("y: ", optimal_y)
	print("Objective value: ", optimal_value)
	
	
	x_values = np.linspace(0, 10, 100)
	y1_values = (12-2*x_values)/3
	y2_values = (18-4*x_values)/2
	
	plt.plot(x_values, y1_values, label=2*x + 3*y >= 12)
	plt.plot(x_values, y2_values, label=4*x + 2*y >= 18)
	
	plt.fill_between(x_values, np.maximum(y1_values,y2_values), 0, where=(y1_values<=y2_values), color="gray")
	plt.fill_between(x_values, np.maximum(y1_values,y2_values), 0, where=(y1_values>=y2_values), color="green", alpha=0.5, label='Unwanted Region')
	
	
	plt.title("Graph showing the feasible region")
	plt.xlim(0,8)
	plt.ylim(0,11)
	plt.grid(True)
	plt.scatter(optimal_x, optimal_y, color="yellow", marker='*', label="Optimal Solution")
	plt.legend()
	plt.show()
\end{verbatim}

\textbf{Solution And Graph}\newline\newline
\includegraphics{solutions/num4}\newline\newline
\includegraphics{C:/Users/skurge/OneDrive/Desktop/graph4}
\newpage




\section{Transportation Logistics 
	Optimization}
\subsection{\textbf{Number Five: Production Planning in Manufacturing}}


\begin{verbatim}
	#libraries
	import numpy as np
	import matplotlib.pyplot as plt
	from mpl_toolkits.mplot3d import Axes3D
	from pulp import LpVariable, LpMinimize, LpProblem
	
	#the problem
	problem = LpProblem(name="Production_minimizing", sense=LpMinimize)
	
	#decision variables
	x1 = LpVariable(name="x", lowBound=0)
	x2 = LpVariable(name="y", lowBound=0)
	x3 = LpVariable(name="z", lowBound=0)
	
	# Define the objective function coefficients
	problem += 5*x1 + 3*x2 + 4*x3, "objective"
	
	# Coefficients of the inequality constraints
	problem += 2*x1 + 3*x2 + x3 <= 1000, "raw_materials"
	problem += 4*x1 + 2*x2 + 5*x3 <= 120, "labour"
	problem += x1 >= 200
	problem += x2 >= 300
	problem += x3 >= 150
	
	
	# Solve 
	problem.solve()
	
	# Display the results
	print("OPTIMAL SOLUTION:")
	print(f"X: {x1.varValue}")
	print(f"Y: {x2.varValue}")
	print(f"Z: {x3.varValue}")
	print(f"Minimum cost: {problem.objective.value()}")
	
	
	# plotting the graph
	
	# Create a meshgrid for x1, x2, and x3
	x1_vals = np.linspace(0, 400, 50)
	x2_vals = np.linspace(0, 400, 50)
	x1_grid, x2_grid = np.meshgrid(x1_vals, x2_vals)
	
	# Calculate the corresponding z-values (objective function)
	z_vals = 5 * x1_grid + 3 * x2_grid + 4 * (1950 - x1_grid - x2_grid)
	
	# Create the 3D plot
	fig = plt.figure(figsize=(10, 7))
	ax = fig.add_subplot(111, projection='3d')
	
	# Plot the feasible region (constraints)
	ax.plot([200, 200], [0, 400], [0, 0], color='yellow', linestyle='--', linewidth=2, label='Raw Materials')
	ax.plot([0, 400], [300, 300], [0, 0], color='blue', linestyle='--', linewidth=2, label='Labor')
	ax.plot([0, 400], [0, 400], [1950, 1950], color='red', linestyle='--', linewidth=2, label='Minimum Production')
	
	
	optimum_x1 = 200
	optimum_x2 = 300
	optimum_z = 1950
	ax.scatter(optimum_x1, optimum_x2, optimum_z, color='purple', s=100, label='Optimum Point')
	
	
	ax.set_xlabel('x1')
	ax.set_ylabel('x2')
	ax.set_zlabel('x3')
	ax.set_title('Production Cost Minimization')
	ax.legend()
	plt.show()
\end{verbatim}
\newpage
\textbf{Solution And Graph}\newline\newline
\includegraphics{solutions/num5}\newline\newline
\includegraphics{C:/Users/skurge/OneDrive/Desktop/graph5}
\newpage




\subsection{\textbf{Number Six: Financial Portfolio Optimization}}


\begin{verbatim}
	#libraries
	import numpy as np
	import matplotlib.pyplot as plt
	from mpl_toolkits.mplot3d import Axes3D
	from pulp import LpVariable, LpMaximize, LpProblem
	
	#the problem
	problem = LpProblem(name="Financial_portfolio_optimization", sense=LpMaximize)
	
	#decision variables
	x1 = LpVariable(name="x", lowBound=0)
	x2 = LpVariable(name="y", lowBound=0)
	x3 = LpVariable(name="z", lowBound=0)
	
	# Define the objective function coefficients
	problem += 0.08*x1 + 0.1*x2 + 0.12*x3, "objective"
	
	# Coefficients of the inequality constraints
	problem += 2*x1 + 3*x2 + x3 <= 10000, "budget"
	problem += x1 >= 2000
	problem += x2 >= 1500
	problem += x3 >= 1000
	
	
	# Solve linear programming problem
	problem.solve()
	
	# Display the results
	print("OPTIMAL SOLUTION:")
	print(f"X: {x1.varValue}")
	print(f"Y: {x2.varValue}")
	print(f"Z: {x3.varValue}")
	print(f"Minimum cost: {problem.objective.value()}")
	
	
	# Create a meshgrid for A, B, and C
	A_vals = np.linspace(0, 2500, 50)
	B_vals = np.linspace(0, 2000, 50)
	A_grid, B_grid = np.meshgrid(A_vals, B_vals)
	
	# Calculate the corresponding z-values (ROI function)
	C_vals = (10000 - 2 * A_grid - 3 * B_grid)  # Budget constraint
	ROI_vals = 0.08 * A_grid + 0.1 * B_grid + 0.12 * C_vals
	
	# Create the 3D plot
	fig = plt.figure(figsize=(10, 6))
	ax = fig.add_subplot(111, projection='3d')
	
	# Plot the feasible region (constraints)
	ax.plot([2000, 2000], [0, 2000], [0, 470], color='green', linestyle='--', linewidth=2, label='Budget')
	ax.plot([0, 2500], [1500, 1500], [0, 470], color='blue', linestyle='--', linewidth=2, label='Maximum Investment')
	ax.plot([0, 2500], [0, 2000], [470, 470], color='red', linestyle='--', linewidth=2, label='Maximum Investment')
	
	# Highlight the optimum point
	optimum_A = 2000
	optimum_B = 1500
	optimum_ROI = 470
	ax.scatter(optimum_A, optimum_B, optimum_ROI, color='yellow', s=100, label='Optimum Point')
	
	# Set labels and title
	ax.set_xlabel('A')
	ax.set_ylabel('B')
	ax.set_zlabel('ROI')
	ax.set_title('Return On Investiment Maximization')
	
	# Add a legend
	ax.legend()
	
	# Show the plot
	plt.show()
\end{verbatim}
\newpage

\textbf{Solution And Graph}\newline\newline
\includegraphics{solutions/num6}\newline\newline
\includegraphics{C:/Users/skurge/OneDrive/Desktop/graph6}

\newpage




\section{\textbf{Project Resource Allocation 
	Optimization}}
\subsection{\textbf{Number Seven: Diet Optimization}}


\begin{verbatim}
	from pulp import*
	import matplotlib.pyplot as plt
	import numpy as np
	
	#lpProblem
	model = LpProblem(name="diet_optimization", sense=LpMinimize)
	
	#LpVariables
	x=LpVariable("x",0)
	y=LpVariable("y",0)
	
	#objective function
	model+=3*x + 2*y
	
	#constraints
	model+= 2*x + y >= 20
	model+=3*x + 2*y>= 25
	
	#solving
	model.solve()
	
	#display results
	x=x.varValue
	y=y.varValue
	optimal_value= model.objective.value()
	
	print("status:",LpStatus[model.status])
	print("x",x)
	print("y",y)
	print("optimal", optimal_value)
	
	
	
	# Define the constraints
	def constraint1(x):
	return 20 - 2*x
	
	def constraint2(x):
	return (25 - 3*x) / 2
	
	# Create x range for plotting
	x_values = np.linspace(0, 15, 400)
	
	# Plot constraints
	plt.plot(x_values, constraint1(x_values), label='2x + y >= 20')
	plt.plot(x_values, constraint2(x_values), label='3x + 2y >= 25')
	
	# Shade unwanted region
	x=[0,0,15]
	y=[12.5,-10,-10]
	plt.fill(x,y,color='gray')
	plt.fill(x,y, color="yellow", label="Feasible Region")
	plt.scatter(x, y, color='red', marker='*', label='Optimal Solution')
	
	
	# Set labels and legend
	plt.xlabel('x')
	plt.ylabel('y')
	plt.title('A graph showing the unwanted region')
	plt.legend()
	
	# Show plot
	plt.grid(True)
	plt.axhline(0, color='black',linewidth=0.5)
	plt.axvline(0, color='black',linewidth=0.5)
	plt.show()
	
\end{verbatim}
\newpage
\textbf{Solution And Graph}\newline\newline
\includegraphics{solutions/num7}\newline\newline
\includegraphics{C:/Users/skurge/OneDrive/Desktop/graph7}
\newpage




\subsection{\textbf{Number Eight: Production Planning}}


\begin{verbatim}
	from pulp import *
	import matplotlib.pyplot as plt
	import numpy as np
	
	model = LpProblem(name = "Production-Planning", sense=LpMaximize)
	
	x1 = LpVariable("x1", 0)
	x2 = LpVariable("x2", 0)
	
	model += 5*x1 + 3*x2
	
	model += 2*x1 + 3*x2 <= 60
	model += 4*x1 + 2*x2 <= 80
	
	model.solve()
	
	optimal_x = x1.varValue
	optimal_y = x2.varValue
	optimal_value = model.objective.value()
	
	print("Optimal Solution: ")
	print("x: ", optimal_x)
	print("y: ", optimal_y)
	print("Objective value: ", optimal_value)
	
	
	#plotting the graph
	x_values = np.linspace(0, 100, 1000)
	y1_values = (60-2*x_values)/3
	y2_values = (80-4*x_values)/2
	
	plt.plot(x_values, y1_values, label=2*x1 + 3*x2 <= 60)
	plt.plot(x_values, y2_values, label=4*x1 + 2*x2 <= 80)
	
	plt.fill_between(x_values, np.maximum(y1_values,y2_values), 0, where=(y1_values>=y2_values), color="black")
	plt.fill_between(x_values, np.maximum(y1_values,y2_values), 0, where=(y1_values<=y2_values), color="gray", alpha=0.5, label='Feasible Region')
	
	
	plt.title("Feasible region")
	plt.xlim(0,38)
	plt.ylim(0,50)
	plt.grid(True)
	plt.scatter(optimal_x, optimal_y, color="yellow", marker='*', label="Optimal Solution")
	plt.legend()
	plt.show()
\end{verbatim}

\textbf{Solution And Graph}\newline\newline
\includegraphics{solutions/num8}\newline\newline



\end{document}